% ================================================================
\section{The Arithmetic Standard Model}
\label{sec:standard-model}
% ================================================================

The Arithmetic Standard Model (\lean{Physics/GaugeTheory/*.lean},
7 files, $\sim$102{,}000 bytes, 76+ theorems, zero sorry, zero custom
axioms) establishes a structural dictionary between the gauge symmetry
groups of particle physics and the arithmetic properties of the first
few primes. Unlike the other physics sections of this paper, which
identify analogies at the level of mathematical structure, the Standard
Model section identifies analogies at the level of \emph{symmetry
groups}: the integers themselves generate a gauge theory.

\subsection{The Three Gauge Sectors}

\subsubsection{U(1): The Electromagnetic Sector (Prime 2)}

The $U(1)$ gauge symmetry (\lean{ArithmeticU1.lean}) is carried by
the unique even prime $p = 2$. The key correspondences (all proved):
\begin{itemize}
  \item \textbf{Charge quantization}: Every integer has a well-defined
    2-adic valuation $v_2(n) \in \NN$, which plays the role of
    electric charge.
  \item \textbf{Charge conservation}: $v_2(mn) = v_2(m) + v_2(n)$---the
    total charge is additive under multiplication (proved in
    \lean{u1\_charge\_additive}).
  \item \textbf{Gauge transformation}: Multiplication by powers of~2
    shifts the charge without changing the odd part---this is the
    $U(1)$ gauge transformation.
  \item \textbf{The Dirac sea}: The even integers $2\ZZ$ form a
    sublattice (the ``Dirac sea'') whose spectral properties differ
    qualitatively from the odd integers (the ``particles'')
    ---cf.~the Dark Sector (\S\ref{sec:dark-sector}).
\end{itemize}

\subsubsection{SU(2): The Weak Sector (Prime 3)}

The $SU(2)$ gauge symmetry (\lean{ArithmeticSU2.lean}) is carried by
the prime $p = 3$. The correspondences:
\begin{itemize}
  \item \textbf{Weak isospin}: The residues modulo~3 form a group
    $\ZZ/3\ZZ$ with three elements: $\{0, 1, 2\}$. The non-zero elements
    $\{1, 2\}$ form a ``doublet'' (the weak isospin pair), while $0$
    is the ``neutral'' state.
  \item \textbf{Parity violation}: The prime 3 breaks the
    left-right symmetry of the integers: among primes $> 3$,
    all are $\equiv 1$ or $2 \pmod{3}$, never $\equiv 0$.
    This is the analogue of parity violation in the weak
    interaction.
  \item \textbf{The CKM matrix}: The Gram matrix restricted to
    residues mod~6 (combining the $U(1)$ and $SU(2)$ sectors)
    produces a mixing matrix between charge eigenstates
    and mass eigenstates.
\end{itemize}

\subsubsection{SU(3): The Color Sector (Primes $\geq 5$)}

The $SU(3)$ color symmetry (\lean{ArithmeticSU3.lean},
26 theorems, zero sorry) is carried by the primes $p \geq 5$.
The correspondences:
\begin{itemize}
  \item \textbf{Color confinement}: No prime $p \geq 5$ is
    highly composite (proved in \lean{confinement\_general}).
    Primes are permanently ``confined'' inside composite
    numbers---free quarks (isolated large primes in the highly
    composite spectrum) do not exist.
  \item \textbf{Asymptotic freedom}: The Gram entry
    $G(p,q) \sim \ln(\gcd(p,q))/\max(p,q)$ decays as the
    primes grow---the color interaction weakens at high energy
    (large primes).
  \item \textbf{The proton (6)}: The number $6 = 2 \cdot 3$ is the
    first perfect number: $\sigma(6) = 12 = 2 \cdot 6$.
    Like the proton, it is perfectly stable---the ``forces''
    (divisor sum) exactly balance (proved in \lean{proton\_stability}).
  \item \textbf{Hadron stability}: $6$ is also the smallest
    number combining $U(1)$ and $SU(2)$ sectors ($6 = 2 \times 3$),
    analogous to the proton being the lightest baryon.
\end{itemize}

\subsection{The Pauli Exclusion Principle}

The Arithmetic Pauli module (\lean{ArithmeticPauli.lean}) proves
the number-theoretic Pauli exclusion principle:

\begin{principle}[Pauli Exclusion $=$ Squarefree Condition]
An integer is \emph{squarefree} if and only if no prime
appears more than once in its factorization:
$\mu(n)^2 = 1 \iff n$ is squarefree. This is exactly the
Pauli exclusion principle: no two identical fermions (prime
factors) can occupy the same state (divide $n$ with multiplicity
$> 1$).

The squarefree integers are the ``Pauli-allowed'' states of the
prime number Fock space. The non-squarefree integers (those with
$\mu(n) = 0$) are ``Pauli-forbidden''---they contain a repeated
prime factor, violating exclusion.
\end{principle}

\subsection{The Complete Assembly}

The Standard Model assembly (\lean{ArithmeticStandardModel.lean})
proves 76+ theorems with zero sorry, establishing the complete
dictionary:

\begin{center}
\small
\begin{tabular}{@{}ll@{}}
\toprule
\textbf{Standard Model Parameter} & \textbf{Arithmetic Analog} \\
\midrule
Gauge group $U(1) \times SU(2) \times SU(3)$ & First primes: 2, 3, ($\geq 5$) \\
Particle masses & $\ln(p)$ for each prime $p$ \\
Coupling constants & Gram entries $G(p,q)$ \\
Mixing angles & Eigenvector components of $G$ \\
Higgs VEV & $G(2,2) = (\ln 2\pi - \gamma)/2 - 1/4$ \\
QCD scale $\Lambda$ & Mertens product $e^{-\gamma}/\ln(N)$ \\
Pauli exclusion & Squarefree condition \\
Confinement & Primes $\neq$ highly composite \\
Proton stability & $6$ is perfect \\
\bottomrule
\end{tabular}
\end{center}

\begin{remark}[Zero Free Parameters]
The Standard Model of particle physics has 19 free parameters
(masses, coupling constants, mixing angles) that must be measured
experimentally. The Arithmetic Standard Model has \textbf{zero}:
every ``parameter'' is a theorem about the integers. The gauge
group is determined by the first three primes. The coupling
constants are Gram matrix entries, themselves integrals of
fractional-part functions. The Higgs VEV is a closed-form
expression involving $\ln 2\pi$ and the Euler--Mascheroni constant.
The integers are the only input; everything else is a theorem.
\end{remark}

\subsection{The 1+1D Dirac Equation}

The Dirac module (\lean{Dirac.lean}) formalizes the connection
between the Cathedral and the Hilbert--P\'olya conjecture through
the 1+1D Dirac equation:
\begin{equation}
  (i\gamma^\mu \partial_\mu - m)\psi = 0
  \label{eq:dirac}
\end{equation}
The key structural identification:
\begin{itemize}
  \item The Nyman--Beurling functions $\{1/(kx)\}$ are the
    \emph{scattered states} of a 1+1D massless Dirac fermion
    on the half-line (Burnol, 1998).
  \item The Liouville function $\lambda(n) = (-1)^{\Omega(n)}$ is
    the \emph{fermion parity operator} $(-1)^F$---it acts as
    chirality, creating the even/odd sector decomposition
    visible in the GPU eigenvalue spectrum.
  \item The Gram matrix implicitly simulates the
    \emph{scattering matrix} of the Dirac fermion, with
    primes acting as scattering potentials.
\end{itemize}

\begin{principle}[The Dirac Sea of Composites]
The composites form the ``Dirac sea''---the filled negative-energy
states. The primes are the ``excitations'' above the sea. The
NB distance $d_N^2 \to 0$ is the statement that the Dirac sea
is \emph{stable}: no excitation can create a net positive energy,
because the scattered waves interfere destructively with perfect
efficiency.
\end{principle}
